\documentclass[12pt,a4paper]{report}
\usepackage[utf8]{inputenc}
\usepackage[T1]{fontenc}
\usepackage{amsmath}
\usepackage{amsfonts}
\usepackage{amssymb}
\usepackage{listings}
\usepackage{xcolor}
\usepackage{hyperref}
\usepackage{graphicx}
\usepackage{tcolorbox}
\usepackage{minted}

% Define colors for syntax highlighting
\definecolor{codegray}{rgb}{0.95,0.95,0.95}
\definecolor{codegreen}{rgb}{0,0.6,0}
\definecolor{codepurple}{rgb}{0.58,0,0.82}
\definecolor{backcolour}{rgb}{0.95,0.95,0.92}

% Python code listing settings
\lstset{
    language=Python,
    backgroundcolor=\color{backcolour},
    commentstyle=\color{codegreen},
    keywordstyle=\color{magenta},
    numberstyle=\tiny\color{gray},
    stringstyle=\color{codepurple},
    basicstyle=\ttfamily\footnotesize,
    breakatwhitespace=false,
    breaklines=true,
    captionpos=b,
    keepspaces=true,
    numbers=left,
    numbersep=5pt,
    showspaces=false,
    showstringspaces=false,
    showtabs=false,
    tabsize=4
}

\title{RF-Modeler: Technical Documentation}
\author{Documentation generated on May 22, 2025}
\date{\today}

\begin{document}

\maketitle
\tableofcontents

\chapter{Introduction}
RF-Modeler is a sophisticated application designed for modeling and visualizing mathematical functions and data, with a particular focus on radio frequency (RF) signal analysis. The application features a modern graphical user interface built with PySide6 and incorporates advanced mathematical modeling capabilities.

\chapter{Project Structure}
The project follows a modular architecture with clear separation of concerns:

\begin{itemize}
    \item \textbf{src/} - Main source code directory
    \begin{itemize}
        \item \textbf{core/} - Core mathematical and plotting functionality
        \item \textbf{gui/} - Graphical user interface components
        \item \textbf{utils/} - Utility functions and helpers
        \item \textbf{assets/} - Application resources and icons
    \end{itemize}
    \item \textbf{App.py} - Main application entry point
    \item \textbf{build\_linux.sh} - Linux build script
    \item \textbf{build\_windows.bat} - Windows build script
    \item \textbf{requirements.txt} - Python dependencies
\end{itemize}

\chapter{Core Module}
\section{Function Plotter}
The FunctionPlotter class in \texttt{src/core/function\_plotter.py} is responsible for mathematical function plotting and evaluation.

\begin{tcolorbox}[title=FunctionPlotter Class]
\begin{lstlisting}[language=Python]
class FunctionPlotter:
    def __init__(self):
        # Initialization with display ranges and calculation parameters
        self.full_range = (-10000, 10000)  # Full calculation range
        self.display_range = (-20, 20)      # Default visible range
        self.num_points = {
            'display': 1500,
            'full': 15000
        }
        self.max_y_bound = 1e6
\end{lstlisting}
\end{tcolorbox}

Key methods:
\begin{itemize}
    \item \textbf{prepare\_function\_string(function\_str)}: Converts mathematical notation to numpy-compatible expressions
    \item \textbf{evaluate\_function(function\_str)}: Evaluates mathematical functions and returns x, y coordinates
    \item \textbf{get\_plot\_limits(y)}: Calculates appropriate y-axis limits for plotting
\end{itemize}

\section{\texttt{src/core/function\_plotter.py}}
\subsection{File Purpose}
This file defines the \texttt{FunctionPlotter} class, which is responsible for preparing and evaluating mathematical functions expressed as strings. It handles the conversion of standard mathematical notation (like \texttt{sin(x)}, \texttt{e\^{}x}, \texttt{pi}) into a format that Python's numerical libraries (like NumPy and SciPy) can understand and compute. It also manages settings related to plot ranges and point density for visualization.

\subsection{Class: \texttt{FunctionPlotter}}
\begin{lstlisting}[language=Python, caption={Excerpt from src/core/function\_plotter.py}, basicstyle=\ttfamily\footnotesize]
class FunctionPlotter:
    def __init__(self):
        # Adjust ranges for better performance and display
        self.full_range = (-10000, 10000)  # Full calculation range
        self.display_range = (-20, 20)      # Default visible range
        self.num_points = {
            'display': 1500,            'full': 15000        }
        self.max_y_bound = 1e6

    def prepare_function_string(self, function_str):
        # ...
        pass

    def evaluate_function(self, function_str):
        # ...
        pass

    def get_plot_limits(self, y):
        # ...
        pass
\end{lstlisting}
\subsubsection{\texttt{\_\_init\_\_(self)}}
\begin{itemize}
    \item \textbf{Purpose}: The constructor initializes a new \texttt{FunctionPlotter} object with default settings for plotting.
    \item \textbf{Parameters}:
        \begin{itemize}
            \item \texttt{self}: A reference to the instance of the class itself.
        \end{itemize}
    \item \textbf{Implementation}:
        \begin{itemize}
            \item \texttt{self.full\_range = (-10000, 10000)}: Sets the full range over which calculations might be performed. This is a tuple representing (min\_x, max\_x).
            \item \texttt{self.display\_range = (-20, 20)}: Sets the default visible range on a plot. This is typically a smaller subset of the \texttt{full\_range}.
            \item \texttt{self.num\_points = \{'display': 1500, 'full': 15000\}}: A dictionary defining the number of points to use when generating data for plotting. More points mean a smoother curve.
                \begin{itemize}
                    \item \texttt{'display'}: Number of points for the visible display range.
                    \item \texttt{'full'}: Number of points for the full calculation range.
                \end{itemize}
            \item \texttt{self.max\_y\_bound = 1e6}: Sets a maximum absolute value for y-coordinates. This is used to prevent extremely large y-values (e.g., from asymptotes) from distorting the plot's scale.
        \end{itemize}
\end{itemize}

\subsubsection{\texttt{prepare\_function\_string(self, function\_str)}}
\begin{itemize}
    \item \textbf{Purpose}: This method takes a mathematical function string (e.g., \texttt{"sin(x) + pi"}) and converts it into a Python-evaluable string. Common mathematical constants and functions are replaced by their NumPy or SciPy equivalents (e.g., \texttt{"np.sin(x) + np.pi"}). It also handles exponentiation.
    \item \textbf{Parameters}:
        \begin{itemize}
            \item \texttt{self}: The instance of the class.
            \item \texttt{function\_str} (str): The mathematical function as a string.
        \end{itemize}
    \item \textbf{Returns}: (str) The processed function string, ready for evaluation by Python's \texttt{eval()} function within a context where NumPy (\texttt{np}) and SciPy (\texttt{scipy}) are available.
    \item \textbf{How it Works}:
        \begin{enumerate}
            \item \textbf{Constants Mapping}: It defines a dictionary \texttt{constants} mapping common mathematical constant names (e.g., \texttt{'pi'}, \texttt{'e'}) to their NumPy representations (\texttt{'np.pi'}, \texttt{'np.e'}).
            \item \textbf{Functions Mapping}: It defines a comprehensive dictionary \texttt{functions} mapping common mathematical function names (e.g., \texttt{'sin'}, \texttt{'cos'}, \texttt{'tan'}, \texttt{'exp'}, \texttt{'log'}, \texttt{'sqrt'}, \texttt{'abs'}, \texttt{'arcsin'}, \texttt{'arccos'}, \texttt{'arctan'}, \texttt{'log10'}, \texttt{'log2'}, \texttt{'sinh'}, \texttt{'cosh'}, \texttt{'tanh'}, \texttt{'ceil'}, \texttt{'floor'}, \texttt{'round'}, \texttt{'erf'} (from \texttt{scipy.special}), \texttt{'gamma'} (from \texttt{scipy.special}), \texttt{'beta'} (from \texttt{scipy.special}), \texttt{'sinc'}) to their NumPy or SciPy counterparts.
            \item \textbf{Replacement}:
                \begin{itemize}
                    \item It iterates through the \texttt{constants} dictionary. For each constant, it uses \texttt{re.sub(rf'\b\{const\}\b', replacement, function\_str)} to replace all occurrences of the constant's name (as a whole word) in \texttt{function\_str} with its NumPy equivalent.
                    \item It then iterates through the \texttt{functions} dictionary. For each function, it uses \texttt{re.sub(rf'\b(?<!np\.)\{func\}\b', replacement, function\_str)} to replace function names (as whole words, and not already prefixed with \texttt{np.}) with their NumPy/SciPy versions. The negative lookbehind \texttt{(?<!np\.)} prevents accidental replacement if a function is already in the desired format (e.g. \texttt{np.sin}).
                \end{itemize}
            \item \textbf{Exponentiation}: It replaces the caret symbol \texttt{\^} (commonly used for exponentiation in mathematics) with Python's exponentiation operator \texttt{**} using \texttt{function\_str.replace('^', '**')}.
        \end{enumerate}
\end{itemize}

\subsubsection{\texttt{evaluate\_function(self, function\_str)}}
\begin{itemize}
    \item \textbf{Purpose}: This method evaluates the processed mathematical function string over a range of x-values to produce corresponding y-values, suitable for plotting. It intelligently generates x-values based on the function type and handles domain restrictions.
    \item \textbf{Parameters}:
        \begin{itemize}
            \item \texttt{self}: The instance of the class.
            \item \texttt{function\_str} (str): The original mathematical function string.
        \end{itemize}
    \item \textbf{Returns}: (tuple) A tuple \texttt{(x\_full, y\_full)} containing NumPy arrays of the x and y coordinates after evaluation, filtering of non-finite values, and clipping.
    \item \textbf{How it Works}:
        \begin{enumerate}
            \item \textbf{Prepare String}: It first calls \texttt{self.prepare\_function\_string(function\_str)} to get the Python-evaluable function string.
            \item \textbf{Identify Special Functions}: It checks for the presence of functions with specific domain requirements (e.g., \texttt{log}, \texttt{sqrt}, \texttt{arcsin}, \texttt{arccos}, \texttt{gamma}, \texttt{beta}).
            \item \textbf{Generate Display X-values}: \texttt{x\_display} points are generated using \texttt{np.linspace} over \texttt{self.display\_range}. These points are then filtered based on the domain requirements of any special functions present (e.g., \texttt{x > 0} for log/sqrt, \texttt{-1 <= x <= 1} for arcsin/arccos).
            \item \textbf{Evaluate Display Y-values}: \texttt{y\_display} is calculated by evaluating the processed function string using \texttt{eval()} with \texttt{x\_display}.
            \item \textbf{Generate Full X-values (\texttt{x\_full})}: This is done adaptively:
                \begin{itemize}
                    \item For functions like \texttt{log} or \texttt{sqrt}, \texttt{x\_full} is generated using \texttt{np.logspace} for regions outside \texttt{x\_display} and concatenated with \texttt{x\_display}, focusing on positive values.
                    \item For \texttt{arcsin} or \texttt{arccos}, \texttt{x\_full} is generated using \texttt{np.linspace} strictly within the \texttt{[-1, 1]} domain.
                    \item For \texttt{gamma} or \texttt{beta}, \texttt{x\_full} uses \texttt{np.logspace} for positive values.
                    \item For other functions, \texttt{x\_full} is generated by concatenating \texttt{np.linspace} for ranges outside \texttt{self.display\_range} with \texttt{x\_display}, covering \texttt{self.full\_range}.
                    \item Duplicate x-values are removed using \texttt{np.unique(x\_full)}.
                \end{itemize}
            \item \textbf{Evaluate Full Y-values}: The processed function is evaluated using \texttt{eval()} with \texttt{x\_full} to get \texttt{y\_full}. The evaluation context includes \texttt{numpy} (as \texttt{np}) and \texttt{scipy} (which makes \texttt{scipy.special} available).
            \item \textbf{Filter and Clip}: Non-finite y-values (NaN, infinity) are removed from \texttt{x\_full} and \texttt{y\_full}. The remaining \texttt{y\_full} values are then clipped to the range \texttt{[-self.max\_y\_bound, self.max\_y\_bound]}.
            \item \textbf{Error Handling}: The entire evaluation is wrapped in a \texttt{try-except} block. If an exception occurs, it attempts to raise a more descriptive \texttt{ValueError} for common domain errors (e.g., "Logarithm error: Cannot compute log of zero or negative values", "Square root error: Cannot compute sqrt of negative values", etc.) or a general error message.
        \end{enumerate}
\end{itemize}

\subsubsection{\texttt{get\_plot\_limits(self, y)}}
\begin{itemize}
    \item \textbf{Purpose}: This method calculates appropriate y-axis limits for a plot based on the calculated y-values, ensuring good visibility of the plotted data.
    \item \textbf{Parameters}:
        \begin{itemize}
            \item \texttt{self}: The instance of the class.
            \item \texttt{y} (np.ndarray): An array of y-values.
        \end{itemize}
    \item \textbf{Returns}: (tuple) A tuple \texttt{(min\_y, max\_y)} representing the suggested y-axis limits for the plot.
    \item \textbf{How it Works}:
        \begin{enumerate}
            \item \textbf{Filter Finite Values}: It first creates \texttt{y\_valid} by selecting only the finite values from the input \texttt{y} array using \texttt{np.isfinite(y)}.
            \item \textbf{Handle Empty Data}: If \texttt{y\_valid} is empty (no finite y-values), it returns a default range of \texttt{(-10, 10)}.
            \item \textbf{Calculate Min/Max}: It finds the minimum and maximum of \texttt{y\_valid}.
            \item \textbf{Adaptive Padding}: Based on the range (\texttt{abs(y\_max - y\_min)}), it applies different padding logic:
                \begin{itemize}
                    \item If range < 0.1: Returns \texttt{(mid - 0.5, mid + 0.5)}.
                    \item If range < 1: Returns \texttt{(mid - 1, mid + 1)}.
                    \item If range < 10: Padding is \texttt{max(1, range * 0.15)}.
                    \item If range > 1000: Padding is \texttt{range * 0.05} (smaller proportional padding for very large ranges).
                    \item Otherwise (medium range): Padding is \texttt{range * 0.1}.
                \end{itemize}
            \item The final limits are \texttt{y\_min - padding} and \texttt{y\_max + padding}.
        \end{enumerate}
\end{itemize}

\chapter{GUI Module}
\section{Main Window}
The MathPlotterApp class in \texttt{src/gui/main\_window.py} serves as the main application window.

\begin{tcolorbox}[title=Main Window Features]
\begin{itemize}
    \item Menu bar with file operations and view controls
    \item Interactive plotting canvas
    \item Status bar with real-time information
    \item Custom navigation toolbar
    \item Integration with side menu
\end{itemize}
\end{tcolorbox}

\subsubsection{\texttt{\_\_init\_\_(self, parent=None)}}
\begin{itemize}
    \item \textbf{Purpose}: Constructor for the \texttt{MathPlotterApp} main window.
    \item \textbf{Parameters}:
        \begin{itemize}
            \item \texttt{parent} (QWidget, optional): Parent widget. Defaults to \texttt{None}.
        \end{itemize}
    \item \textbf{Implementation Details}:
        \begin{enumerate}
            \item \textbf{Initialization}: Calls the superclass (\texttt{QMainWindow}) constructor.
            \item \textbf{Function Plotter}: Initializes an instance of \texttt{FunctionPlotter} from \texttt{src.core.function\_plotter}.
            \item \textbf{Window Setup}: Sets the window title to "RF-Modeler - Function Plotter & Regression Analysis" and applies a custom stylesheet using \texttt{Styles.MAIN\_WINDOW\_STYLE}.
            \item \textbf{Central Widget and Layout}: Creates a central widget and a main horizontal layout (\texttt{QHBoxLayout}).
            \item \textbf{Plotting Area Setup}:
                \begin{itemize}
                    \item Creates a \texttt{Figure} and a \texttt{FigureCanvasQTAgg} for Matplotlib plotting.
                    \item Adds an Axes (\texttt{self.ax}) to the figure for 2D plotting.
                    \item Instantiates \texttt{CustomNavigationToolbar} using the canvas and the main window.
                    \item Creates a vertical layout (\texttt{plot\_layout}) to hold the toolbar and the canvas.
                    \item Adds this \texttt{plot\_layout} to the main horizontal layout.
                \end{itemize}
            \item \textbf{Side Menu}: Initializes and adds \texttt{SideMenu(self.function\_plotter, self)} to the main layout.
            \item \textbf{Status Bar}: Creates and sets up a status bar to display messages.
            \item \textbf{Menu Bar}: Calls \texttt{self.\_create\_menu\_bar()} to populate the menu.
            \item \textbf{Initial Plot}: Calls \texttt{self.plot\_function()} to render an initial empty plot.
            \item \textbf{Window Geometry}: Sets an initial window size (e.g., 1200x800) and centers the window on the screen.
            \item \textbf{Connections}: Connects signals from the side menu to appropriate slots in the main window:
                \begin{itemize}
                    \item \texttt{side\_menu.plot\_requested.connect(self.plot\_function)}
                    \item \texttt{side\_menu.clear\_plot.connect(self.clear\_plot)}
                    \item \texttt{side\_menu.status\_message.connect(self.show\_status\_message)}
                    \item \texttt{side\_menu.toggle\_grid.connect(self.toggle\_grid)}
                    \item \texttt{side\_menu.toggle\_legend.connect(self.toggle\_legend)}
                    \item \texttt{side\_menu.update\_plot\_limits.connect(self.update\_plot\_limits)}
                    \item \texttt{side\_menu.export\_plot.connect(self.export\_plot)}
                    \item \texttt{side\_menu.show\_current\_function.connect(self.show\_current\_function\_dialog)}
                \end{itemize}
        \end{enumerate}
\end{itemize}

\subsubsection{\texttt{\_create\_menu\_bar(self)}}
\begin{itemize}
    \item \textbf{Purpose}: Creates and populates the main menu bar for the application.
    \item \textbf{Implementation Details}:
        \begin{itemize}
            \item \textbf{File Menu}: Creates a "File" menu.
                \begin{itemize}
                    \item \textbf{Export Plot Action}: Adds an "Export Plot..." action (icon: \texttt{save.png}), connects its \texttt{triggered} signal to \texttt{self.export\_plot}.
                    \item \textbf{Exit Action}: Adds an "Exit" action (icon: \texttt{exit.png}), connects its \texttt{triggered} signal to \texttt{self.close}.
                \end{itemize}
            \item \textbf{View Menu}: Creates a "View" menu.
                \begin{itemize}
                    \item \textbf{Toggle Grid Action}: Adds a "Toggle Grid" checkable action, connects its \texttt{toggled} signal to \texttt{self.toggle\_grid}. (Initial state might be based on a setting).
                    \item \textbf{Toggle Legend Action}: Adds a "Toggle Legend" checkable action, connects its \texttt{toggled} signal to \texttt{self.toggle\_legend}. (Initial state might be based on a setting).
                \end{itemize}
            \item \textbf{Help Menu}: Creates a "Help" menu.
                \begin{itemize}
                    \item \textbf{About Action}: Adds an "About" action, connects its \texttt{triggered} signal to \texttt{self.show\_about\_dialog}.
                \end{itemize}
        \end{itemize}
\end{itemize}

\subsubsection{\texttt{plot\_function(self, function\_str=None, plot\_data=None, label=None, color=None, linestyle='-', clear\_previous=False)}}
\begin{itemize}
    \item \textbf{Purpose}: Plots a mathematical function or provided data on the Matplotlib canvas.
    \item \textbf{Parameters}:
        \begin{itemize}
            \item \texttt{function\_str} (str, optional): The mathematical function string to plot.
            \item \texttt{plot\_data} (dict, optional): A dictionary containing pre-calculated data to plot, e.g., \texttt{\{'x': x\_vals, 'y': y\_vals, 'type': 'scatter'/ 'line'\}}. If provided, \texttt{function\_str} is ignored for this specific plot item.
            \item \texttt{label} (str, optional): Label for the plot, used in the legend.
            \item \texttt{color} (str, optional): Color for the plot line/points.
            \item \texttt{linestyle} (str, optional): Linestyle for the plot (e.g., '-', '--', ':'). Default is '-'.
            \item \texttt{clear\_previous} (bool, optional): If True, clears previous plots before drawing the new one. Default is False.
        \end{itemize}
    \item \textbf{Implementation Details}:
        \begin{enumerate}
            \item \textbf{Clear Plot}: If \texttt{clear\_previous} is True, clears the axes (\texttt{self.ax.clear()}).
            \item \textbf{Plotting Logic}:
                \begin{itemize}
                    \item \textbf{If \texttt{plot\_data} is provided}:
                        \begin{itemize}
                            \item Extracts \texttt{x}, \texttt{y}, and \texttt{plot\_type} (defaults to 'line') from \texttt{plot\_data}.
                            \item If \texttt{plot\_type} is 'scatter', uses \texttt{self.ax.scatter()}.
                            \item Else (line plot), uses \texttt{self.ax.plot()}.
                            \item Updates status bar: "Plotting provided data..."
                        \end{itemize}
                    \item \textbf{Else if \texttt{function\_str} is provided and not empty}:
                        \begin{itemize}
                            \item Updates status bar: \texttt{f"Evaluating function: {function\_str}"}.
                            \item Tries to evaluate the function using \texttt{self.function\_plotter.evaluate\_function(function\_str)}.
                            \item If successful, plots \texttt{x} and \texttt{y} using \texttt{self.ax.plot()}.
                            \item Sets y-axis limits using \texttt{self.function\_plotter.get\_plot\_limits(y)}.
                            \item Updates status bar: \texttt{f"Plotted function: {function\_str}"}.
                            \item Handles potential \texttt{ValueError} or other exceptions during evaluation, showing an error message in the status bar and a \texttt{QMessageBox}.
                        \end{itemize}
                    \item \textbf{Else (no function and no data, or empty function string for initial plot)}:
                        \begin{itemize}
                            \item Sets default plot limits (e.g., x-axis from \texttt{self.function\_plotter.display\_range}, y-axis from -10 to 10).
                            \item Updates status bar: "Plot area initialized. Enter a function or load data."
                        \end{itemize}
                \end{itemize}
            \item \textbf{Styling and Legend}: Sets plot title, x-label, y-label. If a legend is enabled (\texttt{self.legend\_enabled}), updates and shows the legend.
            \item \textbf{Grid}: If grid is enabled (\texttt{self.grid\_enabled}), shows the grid.
            \item \textbf{Draw Canvas}: Calls \texttt{self.canvas.draw()} to refresh the plot.
        \end{enumerate}
\end{itemize}

\subsubsection{\texttt{update\_plot\_limits(self, x\_min, x\_max, y\_min, y\_max)}}
\begin{itemize}
    \item \textbf{Purpose}: Updates the x and y axis limits of the plot.
    \item \textbf{Parameters}:
        \begin{itemize}
            \item \texttt{x\_min}, \texttt{x\_max}, \texttt{y\_min}, \texttt{y\_max} (float): New axis limits.
        \end{itemize}
    \item \textbf{Implementation}: Sets \texttt{self.ax.set\_xlim(x\_min, x\_max)} and \texttt{self.ax.set\_ylim(y\_min, y\_max)}, then redraws the canvas.
\end{itemize}

\subsubsection{\texttt{clear\_plot(self)}}
\begin{itemize}
    \item \textbf{Purpose}: Clears all plotted lines and data from the plot area and resets labels/title.
    \item \textbf{Implementation}:
        \begin{enumerate}
            \item Clears the axes: \texttt{self.ax.clear()}.
            \item Resets plot elements: Sets title to "Function Plot", x-label to "X-axis", y-label to "Y-axis".
            \item Reapplies grid and legend if they were enabled.
            \item Sets default axis limits (e.g., x from display range, y from -10 to 10).
            \item Redraws the canvas: \texttt{self.canvas.draw()}.
            \item Updates status bar: "Plot cleared."
        \end{enumerate}
\end{itemize}

\subsubsection{\texttt{toggle\_grid(self, enabled)}}
\begin{itemize}
    \item \textbf{Purpose}: Toggles the visibility of the plot grid.
    \item \textbf{Parameters}:
        \begin{itemize}
            \item \texttt{enabled} (bool): True to show grid, False to hide.
        \end{itemize}
    \item \textbf{Implementation}: Sets \texttt{self.grid\_enabled = enabled}, calls \texttt{self.ax.grid(enabled)}, and redraws the canvas.
\end{itemize}

\subsubsection{\texttt{toggle\_legend(self, enabled)}}
\begin{itemize}
    \item \textbf{Purpose}: Toggles the visibility of the plot legend.
    \item \textbf{Parameters}:
        \begin{itemize}
            \item \texttt{enabled} (bool): True to show legend, False to hide.
        \end{itemize}
    \item \textbf{Implementation}: Sets \texttt{self.legend\_enabled = enabled}. If enabled, calls \texttt{self.ax.legend()} if there are any labeled artists; if disabled, calls \texttt{self.ax.get\_legend().remove()} if a legend exists. Redraws the canvas.
\end{itemize}

\subsubsection{\texttt{export\_plot(self)}}
\begin{itemize}
    \item \textbf{Purpose}: Opens a file dialog to save the current plot as an image file (e.g., PNG, JPG, PDF).
    \item \textbf{Implementation}:
        \begin{enumerate}
            \item Uses \texttt{QFileDialog.getSaveFileName} to get the desired file path and type from the user.
            \item If a file path is provided, saves the figure using \texttt{self.figure.savefig(filePath)}.
            \item Shows a success or error message in the status bar.
        \end{enumerate}
\end{itemize}

\subsubsection{\texttt{show\_status\_message(self, message, timeout=5000)}}
\begin{itemize}
    \item \textbf{Purpose}: Displays a message on the status bar for a specified duration.
    \item \textbf{Parameters}:
        \begin{itemize}
            \item \texttt{message} (str): The message to display.
            \item \texttt{timeout} (int, optional): Duration in milliseconds to show the message. Default is 5000ms.
        \end{itemize}
    \item \textbf{Implementation}: Uses \texttt{self.statusBar().showMessage(message, timeout)}.
\end{itemize}

\subsubsection{\texttt{show\_current\_function\_dialog(self, function\_str)}}
\begin{itemize}
    \item \textbf{Purpose}: Displays a dialog showing the currently processed mathematical function string, potentially with LaTeX rendering.
    \item \textbf{Parameters}:
        \begin{itemize}
            \item \texttt{function\_str} (str): The (potentially processed) function string to display.
        \end{itemize}
    \item \textbf{Implementation}:
        \begin{enumerate}
            \item Creates an instance of \texttt{MathExpressionWidget} (from \texttt{src.gui.show\_function}) with the \texttt{function\_str}.
            \item Executes the dialog using \texttt{dialog.exec()}.
        \end{enumerate}
\end{itemize}

\subsubsection{\texttt{show\_about\_dialog(self)}}
\begin{itemize}
    \item \textbf{Purpose}: Displays an "About" dialog with information about the application.
    \item \textbf{Implementation}: Uses \texttt{QMessageBox.about()} to show a simple dialog with application name, version (if any), and a brief description.
\end{itemize}

\subsubsection{\texttt{closeEvent(self, event)}}
\begin{itemize}
    \item \textbf{Purpose}: Handles the window close event. Can be used to prompt the user for confirmation or perform cleanup.
    \item \textbf{Parameters}:
        \begin{itemize}
            \item \texttt{event} (QCloseEvent): The close event object.
        \end{itemize}
    \item \textbf{Implementation} (Example):
        \begin{enumerate}
            \item Creates a \texttt{QMessageBox} to ask "Are you sure you want to exit?".
            \item If the user clicks "Yes", accepts the event (\texttt{event.accept()}), closing the application.
            \item If the user clicks "No", ignores the event (\texttt{event.ignore()}), keeping the application open.
        \end{enumerate}
\end{itemize}

\subsection{Function: \texttt{main()}}
\begin{itemize}
    \item \textbf{Purpose}: The main entry point function to create and run the \texttt{MathPlotterApp}.
    \item \textbf{Implementation}:
        \begin{enumerate}
            \item Creates a \texttt{QApplication} instance: \texttt{app = QApplication(sys.argv)}.
            \item Creates an instance of the main window: \texttt{main\_window = MathPlotterApp()}.
            \item Shows the main window: \texttt{main\_window.show()}.
            \item Starts the Qt event loop: \texttt{sys.exit(app.exec())}.
        \end{enumerate}
\end{itemize}

\section{\texttt{src/gui/show\_function.py}}
\subsection{File Purpose}
This file provides components for displaying mathematical expressions, potentially rendered using Matplotlib's MathText, in a separate dialog or widget. This is useful for showing the user a nicely formatted version of the function they are working with.

\subsection{Class: \texttt{MathExpressionRenderer(QWidget)}}
\begin{itemize}
    \item \textbf{Purpose}: A custom widget designed to render a mathematical expression string using Matplotlib's MathText capabilities.
    \item \textbf{Inherits from}: \texttt{QWidget}
    \item \textbf{Key Attributes}:
        \begin{itemize}
            \item \texttt{expression} (str): The mathematical expression string to render (e.g., \texttt{r"$y = x^2 + \sin(x)$"}).
        \end{itemize}
\end{itemize}
\subsubsection{\texttt{\_\_init\_\_(self, expression, parent=None)}}
\begin{itemize}
    \item \textbf{Purpose}: Initializes the renderer widget.
    \item \textbf{Parameters}:
        \begin{itemize}
            \item \texttt{expression} (str): The mathematical expression string.
            \item \texttt{parent} (QWidget, optional): Parent widget.
        \end{itemize}
    \item \textbf{Implementation}:
        \begin{enumerate}
            \item Calls superclass constructor.
            \item Stores the \texttt{expression}.
            \item Sets a minimum size for the widget (e.g., 300x100).
        \end{enumerate}
\end{itemize}
\subsubsection{\texttt{paintEvent(self, event)}}
\begin{itemize}
    \item \textbf{Purpose}: Overridden from \texttt{QWidget}. Called whenever the widget needs to be repainted.
    \item \textbf{Parameters}:
        \begin{itemize}
            \item \texttt{event} (QPaintEvent): The paint event object.
        \end{itemize}
    \item \textbf{Implementation}:
        \begin{enumerate}
            \item Creates a \texttt{QPainter} for the widget.
            \item Sets rendering hints for antialiasing.
            \item Clears the widget background.
            \item Uses Matplotlib's MathText parser (\texttt{matplotlib.mathtext.MathTextParser('path')}) to render the \texttt{self.expression}.
            \item The \texttt{parse} method of \texttt{MathTextParser} can output to a \texttt{QImage} or draw directly onto a painter if a suitable interface is provided or wrapped. Alternatively, it might render to an internal buffer and then draw that buffer to the \texttt{QPainter}. The goal is to draw the expression centered within the widget.
            \item Font size and color can be customized.
        \end{enumerate}
\end{itemize}

\subsection{Class: \texttt{MathExpressionWidget(QDialog)}}
\begin{itemize}
    \item \textbf{Purpose}: A dialog box that uses \texttt{MathExpressionRenderer} to display a mathematical expression.
    \item \textbf{Inherits from}: \texttt{QDialog}
\end{itemize}
\subsubsection{\texttt{\_\_init\_\_(self, expression, parent=None)}}
\begin{itemize}
    \item \textbf{Purpose}: Initializes the expression display dialog.
    \item \textbf{Parameters}:
        \begin{itemize}
            \item \texttt{expression} (str): The mathematical expression string to display.
            \item \texttt{parent} (QWidget, optional): Parent widget.
        \end{itemize}
    \item \textbf{Implementation}:
        \begin{enumerate}
            \item Calls superclass constructor.
            \item Sets the window title (e.g., "Current Mathematical Expression").
            \item Creates a \texttt{QVBoxLayout} for the dialog.
            \item Instantiates \texttt{MathExpressionRenderer} with the given \texttt{expression}.
            \item Adds the renderer widget to the layout.
            \item Adds an "OK" button (\texttt{QPushButton}) connected to \texttt{self.accept} (to close the dialog).
            \item Sets the layout for the dialog.
            \item Optionally sets a fixed or minimum size for the dialog.
        \end{enumerate}
\end{itemize}

\section{\texttt{src/gui/side\_menu.py}}

\chapter{Utilities Module}
\section{Data Loader}
The data\_loader.py module provides functions for data import and validation.

Key functions:
\begin{itemize}
    \item \textbf{load\_data\_from\_file(filename)}: Loads and processes CSV data
    \item \textbf{validate\_data(x, y)}: Ensures data meets regression requirements
\end{itemize}

\section{Plot Utilities}
The plot\_utils.py module contains specialized plotting functions.

Key functions:
\begin{itemize}
    \item \textbf{generate\_heatmap\_data(x, y, grid\_size)}: Creates heatmap visualization data
    \item \textbf{plot\_regression\_segments(ax, model, x\_fine, y\_pred\_fine, colors)}: Visualizes regression segments
\end{itemize}

\chapter{Application Entry Point}
The App.py file serves as the main entry point for the application.

\begin{tcolorbox}[title=Application Initialization]
\begin{lstlisting}[language=Python]
from src.gui.main_window import main

if __name__ == "__main__":
    main()
\end{lstlisting}
\end{tcolorbox}

\chapter{Build Scripts}
\section{Linux Build}
The build\_linux.sh script handles:
\begin{itemize}
    \item Environment setup
    \item Dependency installation
    \item Application compilation
    \item Resource bundling
\end{itemize}

\section{Windows Build}
The build\_windows.bat script manages:
\begin{itemize}
    \item Windows-specific environment setup
    \item Dependency resolution
    \item Application packaging
    \item Installer creation
\end{itemize}

\chapter{Source Code Documentation}
\label{cha:source_code_documentation}
This chapter provides a detailed walkthrough of the Python source code that powers the RF-Modeler application. Each section corresponds to a Python file within the \texttt{src} directory, explaining its overall purpose and then delving into the specifics of each class, method, or function it contains.

\section{\texttt{src/\_\_init\_\_.py}}
\subsection{File Purpose}
This file serves as the entry point for the \texttt{src} package. In Python, an \texttt{\_\_init\_\_.py} file in a directory indicates that the directory should be treated as a package or sub-package. This allows modules within the package to be imported using the package's namespace.

The primary role of this specific \texttt{src/\_\_init\_\_.py} file is to provide a high-level description of the RF-Modeler application through its docstring.

\subsection{Content}
\begin{verbatim}
"""
RP-Modeler - A tool for modeling and visualizing mathematical functions and data
"""
\end{verbatim}
\subsubsection{Docstring}
\begin{itemize}
    \item \textbf{Purpose}: The docstring at the beginning of the file, enclosed in triple quotes (\texttt{"""..."""}), provides a brief summary of what the RF-Modeler application is.
    \item \textbf{Usage}: This docstring can be accessed programmatically (e.g., using Python's \texttt{help()} function or by inspecting the package's \texttt{\_\_doc\_\_} attribute) and is also useful for developers to quickly understand the package's main goal.
    \item \textbf{Implementation}: It's a simple string literal. There's no executable code in this file.
\end{itemize}

\section{\texttt{src/core/\_\_init\_\_.py}}
\subsection{File Purpose}
This file marks the \texttt{src/core} directory as a Python sub-package. The "core" sub-package is intended to house the fundamental logic of the application, particularly related to mathematical computations, data modeling, and plotting algorithms. This \texttt{\_\_init\_\_.py} file also makes specific classes from modules within the \texttt{core} sub-package directly accessible when importing from \texttt{src.core}.

\subsection{Content Analysis}
\begin{lstlisting}[language=Python, caption={Content of src/core/\_\_init\_\_.py}, basicstyle=\\ttfamily\\footnotesize]
"""Core functionality for mathematical operations and plotting"""

from .function_plotter import FunctionPlotter
from .regression_model import PiecewiseLogarithmicRegression

__all__ = ['FunctionPlotter', 'PiecewiseLogarithmicRegression']
\end{lstlisting}

\subsubsection{Docstring}
\begin{itemize}
    \item \textbf{Purpose}: The docstring \texttt{"""Core functionality for mathematical operations and plotting"""} describes the general responsibility of the \texttt{core} sub-package.
\end{itemize}

\subsubsection{Imports}
\begin{itemize}
    \item \texttt{from .function\_plotter import FunctionPlotter}
        \begin{itemize}
            \item \textbf{Purpose}: This line imports the \texttt{FunctionPlotter} class from the \texttt{function\_plotter.py} module located in the same directory (the dot \texttt{.} indicates a relative import from the current package).
            \item \textbf{Usage}: This makes \texttt{FunctionPlotter} available as \texttt{src.core.FunctionPlotter}.
        \end{itemize}
    \item \texttt{from .regression\_model import PiecewiseLogarithmicRegression}
        \begin{itemize}
            \item \textbf{Purpose}: This line imports the \texttt{PiecewiseLogarithmicRegression} class from the \texttt{regression\_model.py} module, also in the same directory.
            \item \textbf{Usage}: This makes \texttt{PiecewiseLogarithmicRegression} available as \texttt{src.core.PiecewiseLogarithmicRegression}.
        \end{itemize}
\end{itemize}

\subsubsection{\texttt{\_\_all\_\_} Variable
\begin{itemize}
    \item \texttt{\_\_all\_\_ = ['FunctionPlotter', 'PiecewiseLogarithmicRegression']}
        \begin{itemize}
            \item \textbf{Purpose}: The \texttt{\_\_all\_\_} variable is a list of strings that defines the public interface of the \texttt{src.core} package. When a user writes \texttt{from src.core import *}, only the names listed in \texttt{\_\_all\_\_} will be imported.
            \item \textbf{Usage}: It helps to control which names are exported, preventing accidental exposure of internal variables or modules. In this case, it explicitly states that \texttt{FunctionPlotter} and \texttt{PiecewiseLogarithmicRegression} are the intended public components of this package.
            \item \textbf{Implementation}: It's a list containing the string names of the classes that should be considered public.
        \end{itemize}
\end{itemize}

\section{\texttt{src/core/regression\_model.py}}
\subsection{File Purpose}
This file defines the \texttt{PiecewiseLogarithmicRegression} class, which implements an algorithm for fitting a piecewise logarithmic regression model to a given dataset (x, y). This type of model is useful when the relationship between x and y can be approximated by different logarithmic functions over different segments of the x-range. The class handles finding optimal breakpoints that divide the data into segments and then fits a logarithmic model (\texttt{y = a + b*log(x)}) to each segment.

\subsection{Class: \texttt{PiecewiseLogarithmicRegression}}
This class uses libraries like \texttt{scikit-learn} for metrics (\texttt{r2\_score}), \texttt{scipy.linalg} for solving linear systems (potentially for the regression coefficients), and \texttt{scipy.stats} (likely for confidence intervals or statistical tests).

\begin{lstlisting}[language=Python, caption={Excerpt from src/core/regression\_model.py}, basicstyle=\\ttfamily\\footnotesize]
from sklearn.metrics import r2_score
from scipy.linalg import solve # Not directly used in _solve_logarithmic_regression, SVD is used
from scipy import stats
import numpy as np # Assuming np is imported, as it's used extensively

class PiecewiseLogarithmicRegression:
    def __init__(self, max_breakpoints=5, min_segment_length=3, tolerance=1.00, max_iter=100, verbose=False):
        # ...
        pass
    def _solve_logarithmic_regression(self, X, y):
        # ...
        pass
    # ... other methods ...
    def fit(self, x, y):
        # ...
        pass
    def predict(self, x):
        # ...
        pass
    def get_equations(self):
        # ...
        pass
    def get_universal_function(self):
        # ...
        pass
    def get_confidence_intervals(self, x, confidence=0.95):
        # ...
        pass
\end{lstlisting}

\subsubsection{\texttt{\_\_init\_\_(self, max\_breakpoints=5, min\_segment\_length=3, tolerance=1.00, max\_iter=100, verbose=False)}}
\begin{itemize}
    \item \textbf{Purpose}: Initializes the \texttt{PiecewiseLogarithmicRegression} model with hyperparameters.
    \item \textbf{Parameters}:
        \begin{itemize}
            \item \texttt{max\_breakpoints} (int): Maximum number of breakpoints to consider. Default is 5.
            \item \texttt{min\_segment\_length} (int): Minimum number of data points required in each segment. Default is 3.
            \item \texttt{tolerance} (float): Tolerance for improvement in R-squared to continue adding breakpoints. Default is 1.00 (meaning 100\% R-squared, effectively seeking best fit up to \texttt{max\_breakpoints}).
            \item \texttt{max\_iter} (int): Maximum iterations for the breakpoint optimization algorithm. Default is 100.
            \item \texttt{verbose} (bool): If True, prints progress during fitting. Default is False.
        \end{itemize}
    \item \textbf{Attributes Initialized}:
        \begin{itemize}
            \item \texttt{self.max\_breakpoints}: Stores \texttt{max\_breakpoints}.
            \item \texttt{self.min\_segment\_length}: Stores \texttt{min\_segment\_length}.
            \item \texttt{self.tolerance}: Stores \texttt{tolerance}.
            \item \texttt{self.max\_iter}: Stores \texttt{max\_iter}.
            \item \texttt{self.verbose}: Stores \texttt{verbose}.
            \item \texttt{self.breakpoints\_}: List to store the optimal breakpoints found during fitting. Initialized to an empty list.
            \item \texttt{self.models\_}: List to store the fitted logarithmic regression models (coefficients \texttt{[a, b]}) for each segment. Initialized to an empty list.
            \item \texttt{self.r2\_scores\_}: List to store the R-squared score for each segment. Initialized to an empty list.
            \item \texttt{self.overall\_r2\_score\_}: Overall R-squared score for the piecewise model. Initialized to None.
        \end{itemize}
\end{itemize}

\subsubsection{\texttt{\_solve\_logarithmic\_regression(self, X, y)}}
\begin{itemize}
    \item \textbf{Purpose}: Fits a logarithmic model of the form \texttt{y = a + b * log(x)} to the given data points (X, y).
    \item \textbf{Parameters}:
        \begin{itemize}
            \item \texttt{X} (np.ndarray): Input feature values (independent variable). Must be positive for \texttt{log(X)}.
            \item \texttt{y} (np.ndarray): Target values (dependent variable).
        \end{itemize}
    \item \textbf{Returns}: (tuple) A tuple \texttt{(params, y\_pred, r2)}, where:
        \begin{itemize}
            \item \texttt{params} (np.ndarray or None): The model parameters \texttt{[a, b]} if the fit is successful and \texttt{X} contains positive values. Returns \texttt{None} if \texttt{X} contains non-positive values or if the segment is too short.
            \item \texttt{y\_pred} (np.ndarray or None): The predicted y values using the fitted model. Returns \texttt{None} if \texttt{params} is \texttt{None}.
            \item \texttt{r2} (float or None): The R-squared value for the fit. Returns \texttt{None} if \texttt{params} is \texttt{None}.
        \end{itemize}
    \item \textbf{How it Works}:
        \begin{enumerate}
            \item \textbf{Input Validation}: Checks if \texttt{X} contains any non-positive values or if the length of \texttt{X} is less than \texttt{self.min\_segment\_length}. If so, returns \texttt{(None, None, None)}.
            \item \textbf{Transform X}: Computes \texttt{log\_X = np.log(X)}.
            \item \textbf{Construct Design Matrix}: Creates a design matrix \texttt{A} where the first column is \texttt{np.ones\_like(log\_X)} (for the intercept \texttt{a}) and the second column is \texttt{log\_X} (for the coefficient \texttt{b}).
            \item \textbf{Solve for Parameters}: Uses Singular Value Decomposition (\texttt{np.linalg.svd}) to solve the linear system \texttt{A @ params = y} for \texttt{params = [a, b]}. This is a robust way to perform least squares regression.
            \item \textbf{Predict}: Calculates \texttt{y\_pred = A @ params}.
            \item \textbf{Calculate R-squared}: Computes the R-squared score using \texttt{r2\_score(y, y\_pred)}. If \texttt{y} has no variance (all values are the same), R-squared is set to 1.0 if \texttt{y\_pred} also matches, otherwise 0.0.
        \end{enumerate}
\end{itemize}

\subsubsection{\texttt{fit(self, x, y)}}
\begin{itemize}
    \item \textbf{Purpose}: Fits the piecewise logarithmic regression model to the input data \texttt{x} and \texttt{y}. It iteratively finds the best breakpoints and fits logarithmic models to the segments.
    \item \textbf{Parameters}:
        \begin{itemize}
            \item \texttt{x} (np.ndarray): The independent variable data.
            \item \texttt{y} (np.ndarray): The dependent variable data.
        \end{itemize}
    \item \textbf{Returns}: \texttt{self} (The fitted model instance).
    \item \textbf{How it Works}:
        \begin{enumerate}
            \item \textbf{Data Preparation}: Sorts \texttt{x} and \texttt{y} based on \texttt{x}. Filters out any data points where \texttt{x <= 0} as logarithmic regression is undefined for non-positive \texttt{x}. If no valid data remains, it prints a warning and returns.
            \item \textbf{Initial Fit (No Breakpoints)}: First, it fits a single logarithmic model to the entire dataset using \texttt{\_solve\_logarithmic\_regression}. The R-squared of this initial fit becomes the \texttt{best\_overall\_r2}.
            \item \textbf{Iterative Breakpoint Search}:
                \begin{itemize}
                    \item The algorithm iterates up to \texttt{self.max\_breakpoints} times or \texttt{self.max\_iter} iterations for optimizing breakpoints.
                    \item In each iteration, it tries to add one new breakpoint.
                    \item For each potential new breakpoint position (ensuring \texttt{min\_segment\_length} on both sides), it splits the data and fits logarithmic models to the new segments.
                    \item It calculates a combined R-squared for the current set of breakpoints. The combined R-squared is a weighted average of the R-squared scores of individual segments, weighted by the variance of \texttt{y} within each segment.
                    \item If adding a new breakpoint improves the \texttt{current\_overall\_r2} significantly (the improvement ratio must be greater than \texttt{self.tolerance}), the breakpoint is kept.
                    \item The process of finding the best single additional breakpoint involves iterating through all possible split points in existing segments.
                \end{itemize}
            \item \textbf{Breakpoint Optimization (Refinement)}: After adding breakpoints, it attempts to refine their positions. For each breakpoint, it searches in a local neighborhood (\texttt{search\_radius}, typically 5 data points) for a slightly different position that might yield a better overall R-squared. This is repeated for \texttt{self.max\_iter} or until no further improvement is found.
            \item \textbf{Store Results}: Once the optimal breakpoints are found (or \texttt{max\_breakpoints} is reached), it stores the final breakpoints in \texttt{self.breakpoints\_}, the model parameters for each segment in \texttt{self.models\_}, the R-squared for each segment in \texttt{self.r2\_scores\_}, and the \texttt{self.overall\_r2\_score\_}.
            \item \textbf{Verbose Output}: If \texttt{self.verbose} is True, it prints information about the fitting process, including the number of breakpoints found and the overall R-squared.
        \end{enumerate}
\end{itemize}

\subsubsection{\texttt{predict(self, x)}}
\begin{itemize}
    \item \textbf{Purpose}: Predicts y-values for new x-values using the fitted piecewise logarithmic model.
    \item \textbf{Parameters}:
        \begin{itemize}
            \item \texttt{x} (np.ndarray or list): The x-values for which to make predictions.
        \end{itemize}
    \item \textbf{Returns}: (np.ndarray) The predicted y-values. Returns an array of NaNs if the model has not been fitted.
    \item \textbf{How it Works}:
        \begin{enumerate}
            \item \textbf{Check if Fitted}: If \texttt{self.models\_} is empty (model not fitted), returns an array of NaNs of the same shape as \texttt{x}.
            \item \textbf{Input Conversion}: Converts \texttt{x} to a NumPy array.
            \item \textbf{Initialization}: Creates an output array \texttt{y\_pred} filled with NaNs.
            \item \textbf{Handle Breakpoints}:
                \begin{itemize}
                    \item It defines segments based on \texttt{self.breakpoints\_}. The first segment goes from \texttt{-np.inf} to the first breakpoint, subsequent segments are between breakpoints, and the last segment goes from the last breakpoint to \texttt{np.inf}.
                    \item It iterates through each segment and its corresponding model \texttt{[a, b]} from \texttt{self.models\_}.
                    \item For each segment, it identifies the x-values from the input \texttt{x} that fall within that segment's range (\texttt{segment\_min < x <= segment\_max}).
                    \item It further filters these x-values to include only those \texttt{x > 0} (as \texttt{log(x)} is undefined otherwise).
                    \item For these valid x-values, it calculates \texttt{y\_pred = a + b * np.log(x)}.
                \end{itemize}
        \end{enumerate}
\end{itemize}

\subsubsection{\texttt{get\_equations(self)}}
\begin{itemize}
    \item \textbf{Purpose}: Returns a list of strings, where each string is the equation of the logarithmic model for a segment.
    \item \textbf{Returns}: (list) A list of strings representing the equations (e.g., \texttt{["y = 1.23 + 4.56 * log(x)", ...]}) or an empty list if the model is not fitted.
    \item \textbf{How it Works}:
        \begin{enumerate}
            \item Iterates through \texttt{self.models\_} (which contains \texttt{[a, b]} for each segment).
            \item For each model, formats the equation string as \texttt{f"y = {a:.2f} + {b:.2f} * log(x)"}.
            \item Appends each formatted string to a list.
        \end{enumerate}
\end{itemize}

\subsubsection{\texttt{get\_universal\_function(self)}}
\begin{itemize}
    \item \textbf{Purpose}: Creates and returns a single Python function that can evaluate the entire piecewise model for any given x-value(s).
    \item \textbf{Returns}: (callable) A function that takes an array of x-values and returns an array of predicted y-values based on the fitted piecewise model. Returns \texttt{None} if the model is not fitted.
    \item \textbf{How it Works}:
        \begin{enumerate}
            \item \textbf{Check if Fitted}: If \texttt{self.models\_} is empty, returns \texttt{None}.
            \item \textbf{Define Inner Function}: Defines an inner function, say \texttt{universal\_func(x\_values)}, which will be returned.
            \item \textbf{Inner Function Logic}: This inner function essentially replicates the logic of the \texttt{predict} method:
                \begin{itemize}
                    \item Converts input \texttt{x\_values} to a NumPy array.
                    \item Initializes \texttt{y\_pred} with NaNs.
                    \item Iterates through segments defined by \texttt{self.breakpoints\_} and \texttt{self.models\_}.
                    \item For each segment, applies the corresponding \texttt{a + b * log(x)} formula to the relevant \texttt{x\_values} (where \texttt{x > 0}).
                    \item Returns the populated \texttt{y\_pred} array.
                \end{itemize}
            \item Returns the \texttt{universal\_func}.
        \end{enumerate}
\end{itemize}

\subsubsection{\texttt{get\_confidence\_intervals(self, x, confidence=0.95)}}
\begin{itemize}
    \item \textbf{Purpose}: Calculates the confidence intervals for the predictions of the piecewise logarithmic model at given x-values.
    \item \textbf{Parameters}:
        \begin{itemize}
            \item \texttt{x} (np.ndarray): The x-values for which to calculate confidence intervals.
            \item \texttt{confidence} (float): The confidence level (e.g., 0.95 for 95\% confidence). Default is 0.95.
        \end{itemize}
    \item \textbf{Returns}: (tuple) A tuple \texttt{(lower\_bound, upper\_bound)}, where each is a NumPy array of the same shape as \texttt{x}, containing the lower and upper confidence bounds for the predictions. Returns \texttt{(None, None)} if the model is not fitted or if \texttt{x} is empty.
    \item \textbf{How it Works}:
        \begin{enumerate}
            \item \textbf{Validation}: Checks if the model is fitted and if \texttt{x} is not empty.
            \item \textbf{Initialization}: Creates \texttt{lower\_bound} and \texttt{upper\_bound} arrays filled with NaNs.
            \item \textbf{Iterate Through Segments}:
                \begin{itemize}
                    \item For each segment defined by \texttt{self.breakpoints\_} and its corresponding model in \texttt{self.models\_}:
                        \begin{enumerate}
                            \item Identify the original data points (\texttt{x\_segment}, \texttt{y\_segment}) that were used to fit this segment's model.
                            \item Filter the input \texttt{x} to get values \texttt{x\_eval} that fall within this segment and are positive.
                            \item If no \texttt{x\_eval} or insufficient data in \texttt{x\_segment} (less than 2 points), continue to the next segment.
                            \item Transform \texttt{x\_segment} and \texttt{x\_eval} using \texttt{np.log}.
                            \item Calculate predicted values \texttt{y\_pred\_segment} for \texttt{x\_eval} using the segment's model.
                            \item Calculate the sum of squared errors (SSE) for the segment fit.
                            \item Calculate the standard error of the estimate (SEE) for the segment.
                            \item For each point in \texttt{log\_x\_eval}:
                                \begin{itemize}
                                    \item Calculate the t-statistic for the given confidence level and degrees of freedom (\texttt{len(x\_segment) - 2}).
                                    \item Calculate the margin of error using the SEE, t-statistic, and leverage of the point. The leverage depends on the point's distance from the mean of \texttt{log(x\_segment)} and the sum of squares of \texttt{log(x\_segment)}.
                                    \item The confidence interval for that point is \texttt{y\_pred ± margin\_of\_error}.
                                \end{itemize}
                            \item Store these bounds in the correct positions in \texttt{lower\_bound} and \texttt{upper\_bound}.
                        \end{enumerate}
                \end{itemize}
            \item Returns \texttt{(lower\_bound, upper\_bound)}.
        \end{enumerate}
\end{itemize}

\end{document}
